\section{\bt{Discussion}{讨论}}
\label{sec:limitations}
\bi{\system\ guarantees that a declared use runs on data that satisfies the revision's conditions; business evidence supplies the meaning, and the conditions cover grain, join cardinality, date role, and date-key completeness. Repair restores structure: the uniqueness rule keeps structurally equivalent fixes from being published, and regression against the agent's own SQL accepts a fix only when the change leaves the learning period untouched; a reviewed reference query would automate more repairs, and an ambiguous repair goes to a human with its candidates. The system-level experiments run on the synthetic fixture and on TPC-DS with template-derived definitions, both under injected changes; the end-to-end study uses one LLM with three independent runs per method, and an earlier run with three LLMs, which used the gold SQL as the G8 reference, shows the same pattern. Extending the conditions to value semantics, recognizing definitions in undeclared SQL, and repairs that deduplicate rows or remap keys are the next steps.}
{\system\ 保证声明了修订的使用在满足该修订条件的数据上执行；业务证据提供含义，条件覆盖粒度、连接基数、日期角色与日期键完整性。修复恢复的是结构：唯一性规则防止发布结构上等价的多个修复，以智能体自己的 SQL 为对照的回归测试只在变化不触及学习期时接受修复；经过审核的参照查询能让更多修复自动完成，有歧义的修复带着候选交给人确认。系统层实验在合成数据和模板定义的 TPC-DS 上运行，变化均为注入；端到端实验使用一个大模型、每种方法独立运行 3 次，此前使用 3 个大模型、以标准答案 SQL 作 G8 对照的一轮呈现相同规律。把条件扩展到取值语义、在未声明的 SQL 中识别定义，以及能去重或重映射键的修复，是下一步工作。}
